\section{Preliminaries and Problem Definition}
\label{sec:problem-definition-section}

This section defines the sequential recommendation task and the augmentation policy optimized by RecDataAgent.

\subsection{Sequential Recommendation and Augmentation Policy}
\label{sec:problem-definition}

\textbf{Sequential recommendation.} For domain $\delta$, let $\mathcal{E}_{\delta}$ denote data units, $s_e\in\mathcal{S}_{\delta}$ their train-visible interaction histories and item context, and $D^0_{\delta}$ the immutable base corpus. For a candidate corpus $D$, fixed-protocol training yields parameters $\theta_D$; the recommender predicts $\widehat y_e=f_{\theta_D}(s_e)$, and validation returns aggregate outcomes $\mathbf{y}_D$.

\textbf{Augmentation policy.} Let $\mathcal{R}_{\delta}$ be the catalog of typed data-construction methods. A policy $\pi$ maps the train-visible contexts to user routes, methods, and exposures, inducing an appended corpus $\widehat D_{\delta}(\pi)$ and complete corpus $D_{\delta}(\pi)=D^0_{\delta}\mathbin{\uplus}\widehat D_{\delta}(\pi)$. The downstream learner returns validation outcomes $\mathbf{y}_{\delta}(\pi)=G_{\delta}(D_{\delta}(\pi))$ and validation utility $J_{\delta}(\pi)=\mathcal{U}_{\delta}(\mathbf{y}_{\delta}(\pi))$. Test outcomes are used only for performance reporting and do not enter policy search.

\subsection{Policy Selection}
\label{sec:data-policy-selection}

\textbf{Policy selection.} Before training, the policy observes $\mathcal{E}_{\delta}$, $\{s_e:e\in\mathcal{E}_{\delta}\}$, $D^0_{\delta}$, $\mathcal{R}_{\delta}$, the released archive $\mathcal{H}_{t}$, and fixed exposure budget $Q_{\delta}$. It searches $\pi\in\Pi_{\delta}^{\mathrm{feas}}$ to maximize $J_{\delta}(\pi)$; because the next records and outcomes are unavailable, routing, assignment, and exposure must be searched jointly. RecDataAgent combines train-visible features, typed execution, and aggregate validation feedback for this search.

\section{Method}
\label{sec:method}

Once the policy objective and the information boundary are fixed, the remaining problem is to construct, execute, and revise complete policies under the typed data interfaces.
RecDataAgent separates this procedure into policy formulation, typed execution, archive search, and policy revision.
Figure~\ref{fig:overview} summarizes the resulting policy--adapter interface.

\begin{figure*}[t]
\centering
\includegraphics[width=.99\textwidth]{figures/recdataagent_overview_v31.pdf}
\caption{\method{} as a typed policy-search loop. The controller proposes a complete user-conditioned policy; schema and run validators materialize only feasible records; downstream validation supplies outcome evidence; and the archive uses outcomes and construction diagnostics to form the next candidate batch.}
\label{fig:overview}
\end{figure*}

\subsection{Policy Formulation}

We now define the policy object independently of any particular method implementation. The construction has three stages: encode train-visible context, assign typed actions with explicit exposure, and evaluate the resulting complete corpus through the downstream learner.

For downstream domain $\delta$, let $\mathcal{E}_{\delta}$ be its data units and $s_e\in\mathcal{S}_{\delta}$ the train-visible source context of unit $e\in\mathcal{E}_{\delta}$. 
Let $D^0_{\delta}$ denote the immutable base corpus and let $B_{\delta}(e)=\{b\in D^0_{\delta}:e(b)=e\}$ be the (multi)set of base records owned by $e$.
A deterministic encoder converts this source context into the representation
\begin{equation}
 z_e=\phi_{\delta}(s_e)\in\mathcal{Z}_{\delta}.
 \label{eq:abstract_context}
\end{equation}
The context can cover behavioral, semantic, temporal, and structural information available before downstream evaluation.
Because the context is computed before downstream evaluation, it does not include evaluation labels, predictions, or outcomes.

A typed transformation module serves as a data-construction method.
For a method $o$ with parameter $\eta$, the typed transformation has the form
\begin{equation}
 \begin{aligned}
 o_{\eta}&=\left(A_{o,\eta},T_{o,\eta},\Gamma_{o,\eta}\right),&
 A_{o,\eta}&:\mathcal{Z}_{\delta}\rightarrow\{0,1\},\\
 (T_{o,\eta},\Gamma_{o,\eta})&:
 \mathcal{S}_{\delta}
 \rightarrow\mathsf{View}_{\delta}^{*}\times\mathsf{Trace}_{\delta}^{*}.
 \end{aligned}
 \label{eq:method_tuple}
\end{equation}
The superscript $*$ serves as notation for a finite collection.
Here $A_{o,\eta}$ decides applicability, $T_{o,\eta}$ emits train-visible views, and $\Gamma_{o,\eta}$ returns an auditable trace for each emitted view.
If $A_{o,\eta}(z_e)=0$, no view is emitted and an explicit inapplicable status is recorded.
Methods never receive base-record targets; the materializer binds each view to a base record instead.
The interface is shared across sequence views, retrieved context, relabeled examples, and synthetic views, requiring each to declare its inputs, outputs, applicability, and provenance.

A policy coordinates contextual routing and corpus-level exposure allocation through the following representation:
\begin{equation}
 \pi=\left(\mathbf{r},\nu,\kappa\right),\qquad
 \mathbf{r}=\big((g_1,o_1,\eta_1),\ldots,(g_R,o_R,\eta_R),(g_0,o_0,\eta_0)\big),
 \label{eq:policy_tuple}
\end{equation}
where $g_j:\mathcal{Z}_{\delta}\rightarrow\{0,1\}$ is a route predicate and the final tuple specifies an explicit remainder action.
The allocation rule $\nu$ maps the context population and declared route quotas to a nonnegative exposure vector $\boldsymbol{\nu}_{\pi}\in\mathbb{N}_0^{|\mathcal{E}_{\delta}|}$, where $\nu_{\pi,e}$ is the exposure assigned to unit $e$.
For each unit, first-match routing has the form
\begin{equation}
 a_{\pi}(e)=
 \begin{cases}
 (o_j,\eta_j), & j=\min\{\ell:g_\ell(z_e)=1\},\\
 (o_0,\eta_0), & \text{if no non-remainder predicate fires}.
 \end{cases}
 \label{eq:policy_action}
\end{equation}
Let $T_{a_{\pi}(e)}$ denote $T_{o_j,\eta_j}$ when $a_{\pi}(e)=(o_j,\eta_j)$.
The constraint set $\kappa$ instantiates domain-independent requirements with an exposure budget, coverage, identity preservation, label invariance, and applicability.
The adapter covers any additional constraints specific to the domain.

The policy produces an append multiset and couples it to the immutable base to form the complete training corpus:
\begin{equation}
 \begin{aligned}
 \widehat D_{\delta}(\pi)
 &=\biguplus_{e\in\mathcal{E}_{\delta}}
 \operatorname{Bind}_{\delta}\!\left(
 B_{\delta}(e),
 \mathsf{Expose}_{\nu_{\pi,e}}\!
 \left(T_{a_{\pi}(e)}(s_e)\right)\right),\\
 D_{\delta}(\pi)
 &=D^0_{\delta}\mathbin{\uplus}\widehat D_{\delta}(\pi).
 \end{aligned}
 \label{eq:induced_corpus}
\end{equation}
At exposure $q$, $\mathsf{Expose}_{q}:\mathsf{View}_{\delta}^{*}\to\mathsf{View}_{\delta}^{*}$ is the deterministic rule for selecting or repeating views; no record is emitted when $q=0$.
$\operatorname{Bind}_{\delta}:\mathsf{Base}_{\delta}^{*}\times\mathsf{View}_{\delta}^{*}\to\mathsf{Record}_{\delta}^{*}$ attaches views to base records and leaves the fields declared immutable by the adapter unchanged.
This separation keeps transformation content distinct from supervision quantity and leaves base targets unavailable to method decisions.
Let $D_{\delta}^{\mathrm{src}}=\{(e,s_e,B_{\delta}(e))\}_{e\in\mathcal{E}_{\delta}}$ denote the source package.
Policy feasibility depends on passing the adapter checks for the declarative schema, the declared constraints, and the realized record--trace pair:
\begin{equation}
 \Pi_{\delta}^{\mathrm{feas}}
 =\left\{\pi\in\Pi_{\delta}:
 V_{\delta}^{\mathrm{sch}}(\pi)=1,\ 
 \pi\models\kappa,\ 
 V_{\delta}^{\mathrm{run}}\!\left(\pi,
 X_{\delta}(\pi,D_{\delta}^{\mathrm{src}})\right)=1\right\}.
 \label{eq:feasible_policy}
\end{equation}
The schema check validates the declaration before materialization; the run check validates the record--trace pair returned by $X_{\delta}$ below.
Intended allocation, realized records, and contract-check status are kept distinct throughout execution, preventing one from being mistaken for another.

After corpus construction, the downstream learner supplies the only utility used to evaluate a complete policy:
\begin{equation}
 \mathbf{y}_{\delta}(\pi)
 =G_{\delta}\!\left(D_{\delta}(\pi)\right),\qquad
 J_{\delta}(\pi)=\mathcal{U}_{\delta}\!\left(\mathbf{y}_{\delta}(\pi)\right),\qquad
 \pi_{\delta}^{\star}\in
 \underset{\pi\in\Pi_{\delta}^{\mathrm{feas}}}{\operatorname{arg\,max}}\ J_{\delta}(\pi).
 \label{eq:domain_objective}
\end{equation}
$G_{\delta}$ follows the fixed training, validation, and randomness protocol. Equation~\eqref{eq:domain_objective} defines the search objective; execution selects the highest-validation-utility completed candidate.
This utility cannot be replaced by the planner: a method-local surrogate is not the complete-policy outcome.
The downstream learner evaluates the complete corpus and couples its quality to every local assignment through this black-box utility.
The concrete validation utility, aggregation procedure, and tie-breaking rule used in our experiments are specified in Appendix~\ref{app:experiments}.

Let $\Pi_{\delta}^{\mathrm{contextual}}$ denote the policies above and let $\Pi_{\delta}^{\mathrm{global}}$ be the subset with context-independent action selection and exposure allocation.
A global configuration is the special case with one context-independent action, including a fixed mixture supplied by a deterministic merge action:
\begin{equation}
 \Pi_{\delta}^{\mathrm{global}}\subseteq
 \Pi_{\delta}^{\mathrm{contextual}}.
 \label{eq:policy_containment}
\end{equation}
If a domain contains at least two distinguishable contexts and the method catalog offers compatible actions with different outputs, the inclusion in \eqref{eq:policy_containment} is strict.
Contextual allocation enlarges the policy class and leaves the downstream learner unchanged.

\subsection{Typed Execution}

The framework separates domain-specific adaptation from domain-independent policy search. This is the execution layer: the policy declares what should happen, the adapter determines how a valid record and trace are materialized, and the validator decides whether the realized result can enter the evidence archive.
For domain $\delta$, the adapter interface has the form
\begin{equation}
 \mathcal{I}_{\delta}
 =\left(\phi_{\delta},\mathcal{R}_{\delta},
 V_{\delta}^{\mathrm{sch}},V_{\delta}^{\mathrm{run}},
 X_{\delta},G_{\delta}\right),
 \label{eq:domain_adapter}
\end{equation}
Here $\mathcal{R}_{\delta}$ is the catalog of typed sources and data-construction methods, $V_{\delta}^{\mathrm{sch}}$ validates the declarative recipe, $X_{\delta}$ materializes it deterministically (including $\operatorname{Bind}_{\delta}$), $V_{\delta}^{\mathrm{run}}$ checks realized records and provenance, and $G_{\delta}$ measures the downstream task.
The controller uses these interfaces without either accessing domain-specific files or invoking methods outside the catalog.

With $\Pi_{\delta}^{\mathrm{sch}}=\{\pi\in\Pi_{\delta}:V_{\delta}^{\mathrm{sch}}(\pi)=1\}$, the materializer and run validator have the typed signatures
\begin{equation}
 \begin{aligned}
 X_{\delta}&:
 \Pi_{\delta}^{\mathrm{sch}}\times D_{\delta}^{\mathrm{src}}
 \longrightarrow
 \mathsf{Record}_{\delta}^{*}\times\mathsf{Trace}_{\delta}^{*},\\
 V_{\delta}^{\mathrm{run}}&:
 \Pi_{\delta}^{\mathrm{sch}}\times
 \left(\mathsf{Record}_{\delta}^{*}\times\mathsf{Trace}_{\delta}^{*}\right)
 \longrightarrow\{0,1\}.
 \end{aligned}
 \label{eq:materializer_interface}
\end{equation}
For every emitted record $r$, let $\beta(r)\in D^0_{\delta}$ be its base binding and $\gamma(r)$ its method trace returned with the view.
Let $\mathcal{F}_{\delta}^{\mathrm{fix}}$ denote the fields declared immutable by the task adapter.
The record contract has the form
\begin{equation}
 \begin{aligned}
 e(r)&=e(\beta(r)),&
 f(r)&=f(\beta(r))\quad\forall f\in\mathcal{F}_{\delta}^{\mathrm{fix}},\\
 \operatorname{trace}(r)&=(\beta(r),o,\eta,\gamma(r)).
 \end{aligned}
 \label{eq:record_contract}
\end{equation}
For a schema-valid policy, $X_{\delta}$ converts the source package into the multiset in \eqref{eq:induced_corpus}.
The contract preserves record identity and leaves the task constraints unchanged across implementations.

Routes are evaluated over the context package and combined by a deterministic merge method.
If $R_j(\pi)$ is the first-match set induced by $g_j$,
\begin{equation}
 \displaystyle
 R_j(\pi)=\{e:g_j(z_e)=1\}\setminus\bigcup_{\ell<j}R_\ell(\pi),
 \quad
 R_0(\pi)=\mathcal{E}_{\delta}\setminus\bigcup_{j=1}^{R}R_j(\pi).
 \label{eq:abstract_routes}
\end{equation}
These sets are a disjoint partition before applicability filtering.
Unsupported actions are retained as explicit trace outcomes, preventing their silent removal from the execution record.
For the realized records, exposure and coverage have the form
\begin{equation}
 \displaystyle
 n_e(\pi)=\sum_{r\in\widehat D_{\delta}(\pi)}\mathbf{1}[e(r)=e],
 \quad
  N(\pi)=\sum_{e}n_e(\pi),
 \quad
 \rho(\pi)=\frac{|\{e:n_e(\pi)>0\}|}{|\mathcal{E}_{\delta}|}.
 \label{eq:abstract_accounting}
\end{equation}
The accounting remains valid for multiple emitted views: multiplicity is measured at the record level, and ownership is read from the trace.
Within each policy evaluation, the run validator checks whether realized quantities match allocated exposures and declared corpus constraints.

Here, for RecDataAgent, $\mathcal{E}_{\delta}$ is the user population, $s_e$ supplies train-visible interaction histories and item information, and $\phi_{\delta}$ produces behavioral and semantic features.
Here the catalog contains recommendation data-construction methods, $\mathcal{F}_{\delta}^{\mathrm{fix}}$ contains the instruction and target, $\operatorname{Bind}_{\delta}$ appends each generated view to the input while copying those fields, and $G_{\delta}$ supplies the target recommendation training and evaluation procedure.
These choices instantiate the adapter with recommendation-specific data, records, and evaluation.
The policy grammar, route semantics, materialization contracts, and search controller are unchanged.
Structural portability depends on the separation between the adapter and the search controller.

\paragraph{Structural portability.}
For a new domain $\delta'$, let the adapter supply a context encoder, a typed method catalog, schema and run validators, a deterministic materializer, and an aggregate evaluator satisfying the contracts above.
Replacing $\mathcal{I}_{\delta}$ with $\mathcal{I}_{\delta'}$ keeps the policy grammar, planner, route semantics, archive update, and feedback protocol well-defined, while $\delta'$ determines the feature and record types.
Let $\Pi(\mathcal{R})$ denote the policies formed from catalog $\mathcal{R}$.
Adding a method $\tilde o$ gives the policy-space inclusion
\begin{equation}
 \mathcal{R}_{\delta}'=\mathcal{R}_{\delta}\cup\{\tilde o\}
 \quad\Longrightarrow\quad
 \Pi(\mathcal{R}_{\delta})\subseteq\Pi(\mathcal{R}_{\delta}'),
 \label{eq:method_extension}
\end{equation}
provided that $\tilde o$ satisfies the same typed contract.
These contracts preserve the search interface across domains and methods, with each adapter supplying its task-specific validation objective.
Because no controller rule depends on a particular adapter implementation, substituting the new adapter into every typed call preserves the implication.

\subsection{Archive Search}

Once the policy language and typed execution contract are fixed, the remaining problem is search over complete policies. At round $t$, the search controller maintains the task context and evidence archive and generates candidates according to
\begin{equation}
 C_t=\left(\mathcal{T}_{\delta},\mathcal{F}_{\delta},
 \mathcal{R}_{\delta},\mathcal{H}_t\right),
 \qquad
 \Pi_t=P_{\theta}(C_t)
 =\{\pi_t^{(1)},\ldots,\pi_t^{(K)}\}.
 \label{eq:agent_context}
\end{equation}
Here $\mathcal{T}_{\delta}$ specifies the task and hard constraints, $\mathcal{F}_{\delta}$ summarizes the context population, $\mathcal{R}_{\delta}$ describes method capabilities, and $\mathcal{H}_t$ contains completed validation outcomes and construction diagnostics.
The planner's outcome context consists of validation aggregates, rankings, and incumbent labels; it receives no held-out labels, per-user predictions, or model weights. Source paths are provenance labels, not filesystem access.
In RecDataAgent, a language model implements $P_{\theta}$, while $\operatorname{Update}$ records aggregate validation feedback.
The output is a batch of declarative policies, not executable code.

For each candidate, the typed transition has the form
\begin{equation}
 \pi_t^{(k)}
 \xrightarrow{V_{\delta}^{\mathrm{sch}}}
 (v_t^{(k)},d_t^{(k)})
 \xrightarrow[\ v_t^{(k)}=1\ ]{X_{\delta}}
 (\widehat D_t^{(k)},\tau_t^{(k)})
 \xrightarrow{V_{\delta}^{\mathrm{run}}}
 v_{t,\mathrm{run}}^{(k)}
 \xrightarrow[\ v_{t,\mathrm{run}}^{(k)}=1\ ]{G_{\delta}}
  \mathbf{y}_t^{(k)}
  \xrightarrow{\operatorname{Update}}
 \mathcal{H}_{t+1}.
\label{eq:agent_transition}
\end{equation}

\begin{table*}[!t]
\centering
\caption{Dataset and constructed-corpus statistics for the three Amazon domains. Users, items, interactions, mean history length, sparsity, base records, and appended records are reported for the processed populations and the principal agent-run configurations.}
\label{tab:data-stats}
\small
\setlength{\tabcolsep}{5.5pt}
\begin{tabular}{lrrrrrrr}
\toprule
Domain & Users & Items & Interactions & Avg.\ length & Sparsity & Base records & Added records \\
\midrule
Beauty & 22,363 & 12,101 & 194,687 & 8.71 & 99.9281\% & 34,464 & 22,363 \\
Toys   & 19,412 & 11,924 & 167,597 & 8.63 & 99.9276\% & 31,336 & 19,412 \\
Sports & 35,598 & 18,357 & 296,337 & 8.32 & 99.9547\% & 53,955 & 35,572 \\
\bottomrule
\end{tabular}
\end{table*}

\begin{table*}[!b]
\centering
\caption{Recommendation performance across the three Amazon domains. All baseline methods included in the table retain their method-specific constructions and objectives; \method{} additionally performs user-conditioned data-policy construction and feedback-guided revision. Locally evaluated rows use the matched local protocol, while externally reported rows retain their source protocols. The \method{} row is the validation-selected policy used for final reporting.}
\label{tab:overall}
\small
\setlength{\tabcolsep}{3.2pt}
\renewcommand{\arraystretch}{1.08}
\resizebox{\textwidth}{!}{
\begin{tabular}{lclrrrrrrr}
\toprule
Domain & Group & Method & \multicolumn{2}{c}{Top-5} & \multicolumn{2}{c}{Top-10} & \multicolumn{3}{c}{Top-20} \\
\cmidrule(lr){4-5}\cmidrule(lr){6-7}\cmidrule(lr){8-10}
& & & HR@5 & NDCG@5 & Recall@10 & NDCG@10 & Recall@20 & NDCG@20 & MRR@20 \\
\midrule
\multirow[c]{9}{*}{Beauty} & \multirow[c]{3}{*}{\shortstack[c]{Base\\Models}} & SASRec & 0.0364 & 0.0233 & 0.0550 & 0.0289 & 0.0766 & 0.0338 & 0.0238 \\
& & BERT4Rec & 0.0211 & 0.0133 & 0.0357 & 0.0181 & 0.0493 & 0.0212 & 0.0149 \\
& & HSTU & 0.0388 & 0.0258 & 0.0593 & 0.0322 & 0.0832 & 0.0381 & 0.0262 \\
\cmidrule(lr){2-10}
& \multirow[c]{2}{*}{\shortstack[c]{Generative\\Recommenders}} & OneRec-Think & 0.0507 & 0.0365 & 0.0732 & 0.0441 & 0.0995 & 0.0505 & 0.0361 \\
& & GRLM & 0.0399 & 0.0274 & 0.0591 & 0.0343 & 0.0828 & 0.0403 & 0.0275 \\
\cmidrule(lr){2-10}
& \multirow[c]{3}{*}{\shortstack[c]{Data\\Policy\\Recommenders}} & L2Aug & 0.0468 & 0.0321 & 0.0677 & 0.0390 & 0.0949 & 0.0457 & 0.0319 \\
& & GenPAS & 0.0508 & 0.0351 & 0.0730 & 0.0421 & 0.1019 & 0.0493 & 0.0349 \\
& & AsarRec & 0.0546 & 0.0383 & 0.0784 & 0.0456 & 0.1058 & 0.0531 & 0.0379 \\
\rowcolor{methodblue} & & \method{} &\textbf{0.0563}&\textbf{0.0387}&\textbf{0.0800}&\textbf{0.0463}&\textbf{0.1087}&\textbf{0.0535}&\textbf{0.0380}\\
\midrule
\multirow[c]{9}{*}{Toys} & \multirow[c]{3}{*}{\shortstack[c]{Base\\Models}} & SASRec & 0.0404 & 0.0271 & 0.0570 & 0.0327 & 0.0762 & 0.0375 & 0.0270 \\
& & BERT4Rec & 0.0196 & 0.0120 & 0.0302 & 0.0154 & 0.0408 & 0.0175 & 0.0126 \\
& & HSTU & 0.0333 & 0.0227 & 0.0522 & 0.0281 & 0.0701 & 0.0319 & 0.0230 \\
\cmidrule(lr){2-10}
& \multirow[c]{2}{*}{\shortstack[c]{Generative\\Recommenders}} & OneRec-Think & 0.0527 & 0.0375 & 0.0732 & 0.0441 & 0.0995 & 0.0505 & 0.0361 \\
& & GRLM & 0.0440 & 0.0304 & 0.0626 & 0.0364 & 0.0850 & 0.0420 & 0.0299 \\
\cmidrule(lr){2-10}
& \multirow[c]{3}{*}{\shortstack[c]{Data\\Policy\\Recommenders}} & L2Aug & 0.0484 & 0.0338 & 0.0679 & 0.0402 & 0.0916 & 0.0461 & 0.0333 \\
& & GenPAS & 0.0534 & 0.0368 & 0.0737 & 0.0433 & 0.0987 & 0.0497 & 0.0357 \\
& & AsarRec & 0.0457 & 0.0320 & 0.0654 & 0.0383 & 0.0882 & 0.0441 & 0.0316 \\
\rowcolor{methodblue} & & \method{} &\textbf{0.0592}&\textbf{0.0410}&\textbf{0.0828}&\textbf{0.0487}&\textbf{0.1108}&\textbf{0.0557}&\textbf{0.0401}\\
\midrule
\multirow[c]{9}{*}{Sports} & \multirow[c]{3}{*}{\shortstack[c]{Base\\Models}} & SASRec & 0.0183 & 0.0097 & 0.0272 & 0.0127 & 0.0381 & 0.0147 & 0.0103 \\
& & BERT4Rec & 0.0093 & 0.0059 & 0.0158 & 0.0080 & 0.0218 & 0.0094 & 0.0065 \\
& & HSTU & 0.0245 & 0.0158 & 0.0312 & 0.0209 & 0.0436 & 0.0247 & 0.0167 \\
\cmidrule(lr){2-10}
& \multirow[c]{2}{*}{\shortstack[c]{Generative\\Recommenders}} & OneRec-Think & 0.0264 & 0.0179 & 0.0377 & 0.0219 & 0.0528 & 0.0258 & 0.0176 \\
& & GRLM & 0.0236 & 0.0155 & 0.0376 & 0.0201 & 0.0522 & 0.0235 & 0.0159 \\
\cmidrule(lr){2-10}
& \multirow[c]{3}{*}{\shortstack[c]{Data\\Policy\\Recommenders}} & L2Aug & 0.0245 & 0.0162 & 0.0387 & 0.0207 & 0.0567 & 0.0252 & 0.0166 \\
& & GenPAS & 0.0311 & 0.0204 & 0.0460 & 0.0255 & 0.0668 & 0.0302 & 0.0203 \\
& & AsarRec & 0.0327 & 0.0219 & 0.0481 & 0.0273 & 0.0661 & 0.0319 & 0.0221 \\
\rowcolor{methodblue} & & \method{} &\textbf{0.0385}&\textbf{0.0265}&\textbf{0.0569}&\textbf{0.0324}&\textbf{0.0779}&\textbf{0.0378}&\textbf{0.0261}\\
\bottomrule
\end{tabular}}
\end{table*}


A candidate cannot be materialized before schema validation: its declarative recipe must first pass the schema check.
Here $d_t^{(k)}$ is a schema, execution, or infrastructure diagnostic, $\tau_t^{(k)}$ is a provenance trace, and $\mathbf{y}_t^{(k)}$ is an aggregate validation outcome.
These channels are kept distinct throughout evaluation, preventing any one from being substituted for another.


Let $\mathcal{H}^{+}_t$ denote the complete trials and let $\mathcal{H}^{-}_t$ be the invalid or incomplete attempts.
A complete trial passes both validators and finishes training and validation.
After each trial, its completion status determines the archive update:
\begin{equation}
\begin{aligned}
 \mathcal{H}^{+}_{t+1}
 &=\mathcal{H}^{+}_{t}\cup
 \{(\pi_t^{(k)},\mathbf{y}_t^{(k)},\tau_t^{(k)}):
       \operatorname{complete}_t^{(k)}=1\},\\
 \mathcal{H}^{-}_{t+1}
 &=\mathcal{H}^{-}_{t}\cup
 \{(\pi_t^{(k)},d_t^{(k)}):
       \operatorname{complete}_t^{(k)}=0,\ d_t^{(k)}\neq\varnothing\}.
\end{aligned}
\label{eq:archive_rule}
\end{equation}
$\mathcal{H}^{+}_t$ is the only archive that contributes observations of $J_{\delta}$.
The negative archive preserves actionable failure information without treating an invalid recipe as a low-scoring data point.
The next planner call can retain an evaluated policy, explore another contextual partition, or repair a diagnosed interface mismatch, while leaving the execution semantics unchanged.

\subsection{Policy Revision}

The complete procedure is therefore a typed search-and-evaluation loop rather than an isolated routing rule. The framework's generality depends on this separation of search state from domain-specific execution.
$P_{\theta}$ and $U$ operate on the abstract context and archive, while $\mathcal{I}_{\delta}$ supplies the domain-specific meaning of each feature, method, record, and outcome.
A new dataset, downstream learner, or data-construction capability changes the adapter and objective, while the search state, policy language, and feedback loop remain defined.
The same interface accommodates changes in the domain, method catalog, and downstream model through the corresponding adapter and validation objective.

This ordering yields the experimental progression used in Section~\ref{sec:experiments}: compare complete policies first, intervene on routing and feedback components second, and inspect archive dynamics last.

